\clearpage
\SourceLeaf{458}
\begin{center}Journal für die reine und angewandte Mathematik edited by \textsc{Crelle}.\\
Vol. 22. p. 96.\quad Berlin 1841.\end{center}
\begin{center}\textit{Elementary derivation of a theorem of spherical trigonometry that} \textsc{Legendre} \textit{first stated.}\end{center}
\noindent\hspace*{1.2em}Spherical triangles with small sides may be treated as plane triangles if each spherical angle is diminished by one third of the total spherical excess. The validity of this procedure can be demonstrated in an entirely elementary way as follows.

If we denote the total spherical excess of a spherical triangle by $3\omega$, its three sides by $a,b,c$ and the opposite spherical angles by $A+\omega,B+\omega,C+\omega$, two familiar formulas of spherical trigonometry take the following form:
\[
\begin{aligned}
\sin\tfrac12 a^2&=\frac{\sin\tfrac32\omega\sin(A-\tfrac12\omega)}{\sin(B+\omega)\sin(C+\omega)}\\[5pt]
\cos\tfrac12 a^2&=\frac{\sin(B-\tfrac12\omega)\sin(C-\tfrac12\omega)}{\sin(B+\omega)\sin(C+\omega)},
\end{aligned}
\]
\noindent Combining them gives
\[
\frac{\sin\tfrac12 a^6}{\cos\tfrac12 a^2}
=\frac{\sin\tfrac32\omega^3\sin(A-\tfrac12\omega)^3}
{\sin(B+\omega)^2\sin(B-\tfrac12\omega)\sin(C+\omega)^2\sin(C-\tfrac12\omega)}.
\]
\noindent In the same way we obtain
\[
\frac{\sin\tfrac12 b^6}{\cos\tfrac12 b^2}
=\frac{\sin\tfrac32\omega^3\sin(B-\tfrac12\omega)^3}
{\sin(A+\omega)^2\sin(A-\tfrac12\omega)\sin(C+\omega)^2\sin(C-\tfrac12\omega)}.
\]

